\section{Metodologia da síntese}
Foram selecionados seis fechos de época, todos congelados no commit \texttt{fe11b520\allowbreak bedf1052\allowbreak 57d6f2b2\allowbreak 60ae7b11\allowbreak 0c1a8f61}: encerramento C108--C118, C145, C154, C170, C184 e C189. \evidence{E-S01-EPOCH120B} \evidence{E-S01-C145} \evidence{E-S01-C154} \evidence{E-S01-C170} \evidence{E-S01-C184} \evidence{E-S01-C189}

A síntese privilegia decisões e gates que mudam a direção arquitetural. Não cria um relatório por checkpoint: C119--C137, por exemplo, são herdados apenas quando um fecho posterior os usa como autoridade. Isto reduz repetição sem converter memória editorial em evidência experimental.

A pesquisa bibliográfica tem corte em 30 de setembro de 2026 e privilegia fontes primárias: proceedings/venues oficiais, DOI e arXiv quando necessário. Não é systematic review, meta-analysis ou patent search. Literatura externa sustenta apenas contexto, overlap e claim boundary; não valida números Tesy.

Estados são preservados literalmente quando relevantes: \textbf{MEDIDO}, \textbf{REPRODUZIDO}, \textbf{DIAGNÓSTICO}, \textbf{TRACE\_DERIVED}, \textbf{NOT\_RUN} e \textbf{DESCONHECIDO}. Um protocolo falhado não é reparado editorialmente para PASS.
